\item \textbf{Ciclo de carbono-nitrógeno de Bethe en las estrellas.}
El ciclo del carbono-nitrógeno (ciclo CNO) propuesto por Hans Bethe en 1939 es la fuente principal de energía en estrellas más masivas que el Sol. Identifique analíticamente los núcleos intermedios $A$, $B$, $C$, $D$, $E$ y $F$ formados en las siguientes etapas del ciclo:
\begin{enumerate}
    \item Un protón de alta energía es absorbido por $^{12}\text{C}$ produciendo el núcleo $A$ y un rayo gamma.
    \item El núcleo $A$ decae por emisión de positrones para formar el núcleo $B$.
    \item El núcleo $B$ absorbe un protón para producir el núcleo $C$ y un rayo gamma.
    \item El núcleo $C$ absorbe un protón para producir el núcleo $D$ y un rayo gamma.
    \item El núcleo $D$ decae por emisión de positrones para producir el núcleo $E$.
    \item El núcleo $E$ absorbe un protón para producir el núcleo $F$ y una partícula alfa $^{4}_{2}\text{He}$.
    \item Explique cuál es el significado del núcleo final $F$ al completar la última etapa.
\end{enumerate}